% =====================================================================
% Versão em português do artigo (main.tex é a versão em inglês para a IJCNN).
% No Overleaf: compile este arquivo (Menu > Main document > main_pt.tex).
% Números: historico/teste1_2026-10-16 (relatorio.xlsx, RESULTADOS.md), commit f56b043.
% =====================================================================
\documentclass[conference,a4paper]{IEEEtran}
\IEEEoverridecommandlockouts

\usepackage[utf8]{inputenc}
\usepackage[T1]{fontenc}
\IfFileExists{portuges.ldf}{\usepackage[brazilian]{babel}}{}
\usepackage{cite}
\usepackage{amsmath,amssymb,amsfonts}
\usepackage{graphicx}
\usepackage{booktabs,multirow}
\usepackage{textcomp}
\usepackage{xcolor}
\usepackage[hidelinks]{hyperref}

\newcommand{\todo}[1]{\textcolor{red}{[TODO: #1]}}
\newcommand{\term}[1]{\textsc{#1}}

\begin{document}

\title{Lendo a Rede por Dentro: Regras Fuzzy Aditivas Interpretáveis\\
para Detectar Erros de Classificadores Profundos}

% ---- versão para submissão (duplo-cego) ----
\author{\IEEEauthorblockN{Autores Anônimos}}
% ---- versão final: descomente e apague a linha acima ----
% \author{\IEEEauthorblockN{Matheus Buttow}
% \IEEEauthorblockA{\textit{Programa de Pós-Graduação em ...} \\
% \textit{Universidade ...}\\
% Cidade, Brasil \\
% matheusbuttow@gmail.com}}

\maketitle

\begin{abstract}
Classificadores profundos frequentemente erram com alta confiança, o que é crítico em aplicações médicas como a triagem de câncer de pele,
em que um erro não detectado atrasa o diagnóstico.
Propomos um detector interpretável de erros de classificação que lê o comportamento interno de uma rede já treinada e congelada:
sondas lineares em cada bloco revelam se a rede mudou de ideia no caminho até a resposta,
vistas aumentadas revelam se a resposta é estável, e vizinhos mais próximos no espaço de características revelam se a entrada é familiar.
Esses sinais são combinados por um modelo fuzzy aditivo, um sistema Takagi--Sugeno de ordem zero com regras de um único antecedente e saída
logística, de modo que cada regra tem uma forma legível (``SE a concordância no último bloco é \term{baixa} ENTÃO o risco de erro aumenta'')
e um peso aprendido.
No HAM10000 (benigno vs.\ maligno, EfficientNet-B0) e em seis cenários adicionais (CIFAR-10 com e sem 30\% de ruído nos rótulos,
CIFAR-100, SVHN e MNIST), todos avaliados nas mesmas redes congeladas e nas mesmas imagens de teste que doze detectores concorrentes,
o método obtém a maior AUROC média, ou empata com a maior, em sete de oito cenários e fica a 0,004 do melhor no restante
(AUROC de 0,916 no HAM10000 e de 0,903 no CIFAR-10 com ruído, contra 0,865 e 0,797 da confiança do softmax).
Em relação ao meta-modelo baseado em sondas de Chen et al., é significativamente melhor em 18 de 26 rodadas e nunca significativamente pior,
e suas regras permanecem estáveis entre reamostragens.
\end{abstract}

\begin{IEEEkeywords}
estimação de incerteza, detecção de erros de classificação, sistemas fuzzy, interpretabilidade, classificação de lesões de pele
\end{IEEEkeywords}

% =====================================================================
\section{Introdução}

Redes profundas para dermatoscopia já alcançam acurácia de especialista em bases de referência, mas ainda erram, e alguns desses erros
vêm acompanhados de alta confiança. Em um fluxo de triagem, o custo é assimétrico: uma lesão maligna classificada como benigna com
pontuação confiante dificilmente será revisada. Uma forma prática de usar esse modelo é a \emph{predição seletiva}: a rede responde
os casos em que provavelmente acerta e encaminha os demais a um dermatologista. Para isso, é preciso uma pontuação que ordene as imagens
pelo risco de erro, tarefa conhecida como detecção de erros de classificação (ou detecção de falhas)~\cite{hendrycks2017,jaeger2023}.

A pontuação mais simples é a própria confiança da rede, a maior probabilidade do softmax~\cite{hendrycks2017}. Ela é uma referência forte,
mas é mal calibrada~\cite{guo2017} e, por construção, enxerga apenas a última camada. Aproximações bayesianas como o
MC Dropout~\cite{gal2016} e o aumento de dados no teste ainda resumem a saída, e pontuações baseadas em distância, como
Mahalanobis~\cite{lee2018} e Deep $k$NN~\cite{papernot2018}, visam entradas fora da distribuição, e não erros dentro dela.
Meta-modelos que leem as representações internas, como a abordagem de sondas lineares de Chen et al.~\cite{chen2019},
os vizinhos camada a camada~\cite{font2025} ou classificadores sobre a penúltima camada~\cite{gorbett2022}, são mais informativos,
mas também são caixas-pretas: dizem \emph{que} uma imagem é arriscada, não \emph{por quê}.

Argumentamos que os motivos pelos quais uma rede falha são visíveis na forma como ela chega à resposta, e que esses motivos podem ser
escritos como regras. Uma predição correta e confiante costuma se formar cedo e se manter até o fim; uma errada tende a hesitar entre
os blocos, a mudar quando a imagem é espelhada, ou a ficar entre vizinhos de treino que receberam outra resposta.
Medimos esses comportamentos com um pequeno conjunto de características e os combinamos com um modelo fuzzy aditivo cujas regras um
médico consegue ler.

Nossas contribuições são:
\begin{itemize}
  \item uma leitura da rede congelada, sem retreino, em quatro eixos complementares (a resposta, o caminho pelos blocos,
        a estabilidade entre vistas e a familiaridade da entrada), resumida em poucas características interpretáveis;
  \item um detector fuzzy aditivo, um sistema Takagi--Sugeno de ordem zero com regras de um único antecedente e saída logística, cujos
        pesos são estáveis entre reamostragens (correlação de Spearman de 0,91 a 0,94 no HAM10000);
  \item uma comparação controlada com doze detectores em oito cenários, usando as mesmas redes congeladas, as mesmas imagens de teste e as
        mesmas partições de validação cruzada, com testes de bootstrap pareado, além de três verificações com incerteza real
        (lesões confirmadas por biópsia, várias fotografias da mesma lesão e discordância humana no CIFAR-10H~\cite{peterson2019}).
\end{itemize}

% =====================================================================
\section{Trabalhos Relacionados}

\subsection{Detecção de erros e de falhas}
A maior probabilidade do softmax (MSP)~\cite{hendrycks2017} continua sendo a referência, e um estudo amplo mostrou que a maioria dos
métodos dedicados não a supera de forma consistente para erros dentro da distribuição~\cite{jaeger2023}. O MC Dropout~\cite{gal2016} e os
\textit{ensembles}~\cite{lakshminarayanan2017} estimam a incerteza epistêmica a um custo maior de inferência. O ConfidNet~\cite{corbiere2019}
aprende a prever a probabilidade da classe verdadeira a partir da penúltima camada. As pontuações baseadas em distância, Mahalanobis~\cite{lee2018}
e Deep $k$NN~\cite{papernot2018}, medem a familiaridade no espaço de características e foram pensadas principalmente para detectar
entradas fora da distribuição.

\subsection{Meta-modelos sobre representações internas}
Sondas lineares~\cite{alain2016} mostram o quanto o rótulo é linearmente decodificável em cada camada. Chen et al.~\cite{chen2019} treinam uma
sonda por camada e alimentam um meta-modelo de \textit{gradient boosting} com as saídas de todas as sondas para prever se a rede base acerta;
seus ganhos são maiores sob ruído nos rótulos. Font et al.~\cite{font2025} comparam os vizinhos de treino mais próximos ao longo das camadas,
e Gorbett e Blanchard~\cite{gorbett2022} treinam um classificador sobre a penúltima camada. A profundidade em que uma predição se forma também
já foi relacionada à dificuldade do exemplo~\cite{baldock2021}. Esses trabalhos confirmam que o sinal interno existe, mas seus meta-modelos
não são interpretáveis.

\subsection{Sistemas fuzzy para incerteza e explicação}
Sistemas fuzzy baseados em regras já foram usados para explicar redes profundas a partir de valores de atribuição~\cite{aghaeipoor2023} e para
modelar a incerteza dentro da rede, com camadas fuzzy tipo 2 ou neuro-fuzzy. O aprendizado de regras a partir de dados remonta a Wang e
Mendel~\cite{wangmendel1992}, e os modelos Takagi--Sugeno~\cite{takagi1985} combinam antecedentes linguísticos com consequentes numéricos.
Até onde sabemos, nenhum trabalho anterior aplica regras fuzzy ao comportamento interno de uma rede treinada para detectar seus erros.

% =====================================================================
\section{Método}

Seja $f$ um classificador treinado com $K$ classes, mantido congelado. Para uma entrada $x$, $f$ prevê $\hat y(x)$ com confiança $p(x)$:
no caso binário (HAM10000), $p$ é a saída sigmoide $P(\text{maligno}\mid x)$ e $\hat y = \mathbf{1}[p \ge \tau]$, com o limiar $\tau$
escolhido em um conjunto de validação (índice de Youden); no caso multiclasse, $p$ é a maior probabilidade do softmax.
O objetivo é uma pontuação $r(x)$ que seja alta quando $\hat y(x)$ estiver errada.

\subsection{Lendo a rede}
\label{sec:leitura}

\textit{O caminho pelos blocos.}
Seja $a_b(x)$ a saída do bloco $b \in \{1,\dots,B\}$, com média sobre as posições espaciais. Em um conjunto de imagens de referência do treino,
treinamos uma sonda por bloco, $g_b$, uma regressão logística (com padronização e PCA para 256 dimensões nas camadas largas) cujo alvo é a
\emph{decisão da rede} $\hat y$, e não o rótulo. A sonda responde, portanto: ``se a rede parasse no bloco $b$, o que responderia?''.
Para uma imagem de teste, seja $c_b(x) = g_b(\hat y(x) \mid a_b(x))$ a probabilidade de o bloco $b$ já dar a resposta final, e
$s_b(x) = \mathbf{1}[\arg\max g_b(a_b(x)) = \hat y(x)]$ a sua concordância. Calculamos
\begin{align}
  \textit{trocas} &= \textstyle\sum_{b=2}^{B} |s_b - s_{b-1}|, \\
  \textit{profundidade da decisão} &= \tfrac{1}{B}\max\{b : s_b = 0\}, \\
  \textit{entropia das sondas} &= \tfrac{1}{B}\textstyle\sum_{b} H\big(g_b(a_b)\big)/\log K, \\
  \textit{concordância final} &= \tfrac{1}{|L|}\textstyle\sum_{b \in L} c_b, \\
  \textit{salto final} &= |p(x) - c_B(x)|,
\end{align}
em que $L$ contém os últimos $\lfloor B/3 \rfloor$ blocos; a profundidade da decisão é zero quando todos os blocos concordam com $\hat y$ e
segue a profundidade de predição de Baldock et al.~\cite{baldock2021}.

\textit{Estabilidade entre vistas.}
Cada imagem também passa por $V-1$ transformações que preservam o rótulo (espelhamentos e rotação de 180$^\circ$ para lesões de pele;
espelhamento horizontal e deslocamentos de 2 pixels para o CIFAR; apenas pequenos deslocamentos para SVHN e MNIST, em que espelhar muda o dígito).
Registramos a fração de vistas cuja decisão difere de $\hat y$ (\textit{vistas que mudam}), o desvio padrão da confiança em $\hat y$ entre as
vistas (\textit{desvio entre vistas}) e o desvio padrão do caminho das sondas $(c_1,\dots,c_B)$ entre as vistas, com média sobre os blocos
(\textit{desvio do caminho}).

\textit{Familiaridade.}
No espaço da penúltima camada, recuperamos as $k=15$ imagens de referência mais próximas (similaridade do cosseno) e registramos a fração cuja
decisão da rede difere de $\hat y$ (\textit{discordância dos vizinhos}) e a distância média a elas dividida pela distância mediana entre as
imagens de referência (\textit{distância aos vizinhos}).

Junto com o próprio $p$, essas onze características descrevem a resposta, o caminho, a estabilidade e a familiaridade de cada predição.
As três características de vistas exigem $V$ passagens pela rede; todas as outras exigem uma.

\subsection{Modelo fuzzy aditivo}
Cada característica selecionada $x_j$ é coberta por três funções de pertinência triangulares, \term{baixa}, \term{média} e \term{alta}, que
formam uma partição forte com picos nos percentis 5, 50 e 95. Para a saída binária $p$, os picos são fixos em $\{0, \tau, 1\}$, de modo que
\term{média} significa ``no limiar de decisão''. O risco de erro é
\begin{equation}
  r(\mathbf{x}) = \sigma\Big(b + \sum_{j=1}^{d}\sum_{k=1}^{3} w_{jk}\,\mu_{jk}(x_j)\Big),
  \label{eq:aditivo}
\end{equation}
em que $\mu_{jk}$ é a pertinência de $x_j$ no termo $k$ e $\sigma$ é a função logística. Cada par $(j,k)$ é uma regra
``SE $x_j$ é $k$ ENTÃO o log da chance de erro muda em $w_{jk}$'', de modo que o modelo é um sistema Takagi--Sugeno de ordem zero~\cite{takagi1985}
com regras de um único antecedente cujos consequentes são somados e passados por $\sigma$. Os pesos são aprendidos por regressão logística com
regularização $\ell_2$ sobre as pertinências, tendo como alvo o erro da rede ($\hat y \ne y$). Com $d=4$ características, o modelo tem
12 regras e um viés; um peso positivo significa que a condição aumenta o risco.

\subsection{Seleção de características}
Incluímos sempre $p$, pois a própria resposta é o sinal isolado mais forte, e acrescentamos as $d-1=3$ características com a maior
\emph{AUROC além de $P$}: a AUROC para detectar erros calculada dentro de cinco faixas de quantis de $p$ e ponderada pelo tamanho das faixas.
Ela mede o que uma característica acrescenta entre imagens com confiança parecida. Uma candidata é descartada se sua correlação de Spearman,
em valor absoluto, com uma característica já escolhida for de pelo menos 0,85. A seleção é refeita dentro de cada partição de treino.
Também relatamos uma variante sem as três características de vistas (uma única passagem pela rede).

% =====================================================================
\section{Configuração Experimental}

\subsection{Cenários e redes}
A Tabela~\ref{tab:redes} lista os oito cenários. O HAM10000~\cite{tschandl2018} é dividido por lesão (uma imagem por lesão, de modo que
nenhuma lesão aparece em dois conjuntos) em treino, validação e teste (70/10/20\%), com \textit{mel}, \textit{bcc} e \textit{akiec} como
malignas. Usamos todas as lesões (cinco sementes) e um subconjunto de 3.000 lesões (três sementes). Uma EfficientNet-B0 pré-treinada no
ImageNet, com uma cabeça de duas camadas, é ajustada por 15 épocas (AdamW), e a época e o limiar são escolhidos na validação; seus 18
sub-blocos são lidos pelas sondas.
Os demais cenários reproduzem os dos artigos com que comparamos: ResNet-32 no CIFAR-10, limpo e com 30\% dos rótulos de treino trocados
por outra classe~\cite{chen2019}; VGG-16 no CIFAR-10 e no CIFAR-100, uma CNN pequena no SVHN e uma MLP no MNIST~\cite{corbiere2019,jaeger2023}.
Cada rede é treinada uma vez por semente e mantida congelada; 5.000 imagens de treino servem de referência para as sondas, os vizinhos e o
Mahalanobis.

\subsection{Detectores comparados}
Todos os detectores pontuam as mesmas imagens de teste da mesma rede congelada (pontuação maior = erro mais provável).
\emph{Sem treino:} confiança do softmax (MSP) e entropia~\cite{hendrycks2017}, distância ao limiar (binário), desvio entre vistas (TTA),
MC Dropout com $T=30$ passagens pela cabeça~\cite{gal2016}, Mahalanobis com covariância compartilhada~\cite{lee2018} e
Deep $k$NN~\cite{papernot2018}.
\emph{Com treino:} um modelo só com a saída (regressão logística sobre $p$ e $p^2$ no caso binário; igual ao MSP no multiclasse),
Chen et al.~\cite{chen2019} (sondas treinadas com o rótulo em todos os blocos, as 3 maiores probabilidades de cada sonda, \textit{gradient boosting}),
Font et al.~\cite{font2025} (vizinhos mais próximos camada a camada), Gorbett e Blanchard~\cite{gorbett2022} (SVM linear sobre a penúltima camada)
e uma versão pós-treino do ConfidNet~\cite{corbiere2019} (uma MLP que prevê a probabilidade da classe verdadeira a partir da penúltima camada).
Todos os detectores com treino, inclusive o nosso, usam a mesma validação cruzada estratificada de 5 partições no conjunto de teste e são
avaliados nas pontuações fora da partição.

\subsection{Métricas}
A métrica principal é a AUROC para separar os erros da rede de seus acertos. Também relatamos a AUPR dos erros, a FPR com 95\% de TPR,
o excesso de área sob a curva risco--cobertura (E-AURC)~\cite{jaeger2023} e o número de erros encontrados quando 10\%, 20\% e 30\% das
imagens mais incertas são encaminhadas a um especialista. As diferenças de AUROC são testadas em cada rodada com bootstrap pareado
(1.000 reamostragens comuns a todos os métodos); uma diferença é significativa quando seu intervalo de confiança de 95\% não contém o zero.

% =====================================================================
\section{Resultados}

% Gerado por tabelas_latex.py a partir de relatorio.xlsx.
% Preâmbulo: \usepackage{booktabs,multirow,graphicx}

% ==================================================================
% Tabela 1: Redes
% ==================================================================
\begin{table}[t]
\centering
\caption{Redes avaliadas (congeladas). Média entre sementes.}
\label{tab:redes}
\resizebox{\columnwidth}{!}{%
\begin{tabular}{llcccc}
  \toprule
  Cenário & Rede & Classes & \begin{tabular}[c]{@{}c@{}}Imagens de\\teste\end{tabular} & Acurácia & Sementes \\
  \midrule
  HAM10000 (todas) & EfficientNet-B0 & 2 & 1.495 & 0.861 & 5 \\
  HAM10000 (3000) & EfficientNet-B0 & 2 & 601 & 0.831 & 3 \\
  CIFAR-10 ruído 30\% & ResNet-32 & 10 & 10.000 & 0.865 & 3 \\
  CIFAR-100 & VGG-16 & 100 & 10.000 & 0.731 & 3 \\
  CIFAR-10 (ResNet) & ResNet-32 & 10 & 10.000 & 0.929 & 3 \\
  CIFAR-10 (VGG) & VGG-16 & 10 & 10.000 & 0.937 & 3 \\
  SVHN & CNN & 10 & 26.032 & 0.954 & 3 \\
  MNIST & MLP & 10 & 10.000 & 0.991 & 3 \\
  \bottomrule
\end{tabular}}
\end{table}

% ==================================================================
% Tabela 2: AUROC
% ==================================================================
\begin{table*}[t]
\centering
\caption{AUROC para detectar os erros da rede (média $\pm$ desvio entre sementes). Em negrito, a maior média de cada cenário.}
\label{tab:auroc}
\resizebox{\textwidth}{!}{%
\begin{tabular}{lcccccccc}
  \toprule
  Método & \begin{tabular}[c]{@{}c@{}}HAM10000\\(todas)\end{tabular} & \begin{tabular}[c]{@{}c@{}}HAM10000\\(3000)\end{tabular} & \begin{tabular}[c]{@{}c@{}}CIFAR-10 ruído\\30\%\end{tabular} & CIFAR-100 & \begin{tabular}[c]{@{}c@{}}CIFAR-10\\(ResNet)\end{tabular} & \begin{tabular}[c]{@{}c@{}}CIFAR-10\\(VGG)\end{tabular} & SVHN & MNIST \\
  \midrule
  Softmax (MSP) & 0.865\,{\scriptsize$\pm$0.038} & 0.819\,{\scriptsize$\pm$0.099} & 0.797\,{\scriptsize$\pm$0.014} & 0.866\,{\scriptsize$\pm$0.003} & 0.928\,{\scriptsize$\pm$0.004} & 0.915\,{\scriptsize$\pm$0.004} & \textbf{0.926}\,{\scriptsize$\pm$0.001} & 0.985\,{\scriptsize$\pm$0.000} \\
  Entropia do softmax & 0.865\,{\scriptsize$\pm$0.038} & 0.819\,{\scriptsize$\pm$0.099} & 0.778\,{\scriptsize$\pm$0.015} & 0.867\,{\scriptsize$\pm$0.003} & 0.928\,{\scriptsize$\pm$0.004} & 0.914\,{\scriptsize$\pm$0.004} & 0.924\,{\scriptsize$\pm$0.001} & 0.984\,{\scriptsize$\pm$0.000} \\
  Distância ao limiar & 0.862\,{\scriptsize$\pm$0.013} & 0.842\,{\scriptsize$\pm$0.069} & -- & -- & -- & -- & -- & -- \\
  TTA & 0.864\,{\scriptsize$\pm$0.019} & 0.804\,{\scriptsize$\pm$0.067} & 0.810\,{\scriptsize$\pm$0.012} & 0.845\,{\scriptsize$\pm$0.002} & 0.931\,{\scriptsize$\pm$0.003} & 0.931\,{\scriptsize$\pm$0.005} & 0.909\,{\scriptsize$\pm$0.003} & \textbf{0.986}\,{\scriptsize$\pm$0.000} \\
  MC Dropout & 0.864\,{\scriptsize$\pm$0.039} & 0.818\,{\scriptsize$\pm$0.100} & 0.779\,{\scriptsize$\pm$0.013} & 0.864\,{\scriptsize$\pm$0.002} & 0.926\,{\scriptsize$\pm$0.004} & 0.910\,{\scriptsize$\pm$0.007} & 0.918\,{\scriptsize$\pm$0.002} & 0.984\,{\scriptsize$\pm$0.001} \\
  Mahalanobis & 0.178\,{\scriptsize$\pm$0.028} & 0.214\,{\scriptsize$\pm$0.091} & 0.695\,{\scriptsize$\pm$0.014} & 0.847\,{\scriptsize$\pm$0.002} & 0.802\,{\scriptsize$\pm$0.009} & 0.918\,{\scriptsize$\pm$0.002} & 0.492\,{\scriptsize$\pm$0.009} & 0.769\,{\scriptsize$\pm$0.015} \\
  Deep kNN & 0.610\,{\scriptsize$\pm$0.012} & 0.585\,{\scriptsize$\pm$0.070} & 0.699\,{\scriptsize$\pm$0.008} & 0.667\,{\scriptsize$\pm$0.002} & 0.738\,{\scriptsize$\pm$0.004} & 0.594\,{\scriptsize$\pm$0.013} & 0.859\,{\scriptsize$\pm$0.005} & 0.888\,{\scriptsize$\pm$0.007} \\
  \midrule
  Só P$^\dagger$ & 0.907\,{\scriptsize$\pm$0.017} & 0.872\,{\scriptsize$\pm$0.060} & 0.797\,{\scriptsize$\pm$0.014} & 0.866\,{\scriptsize$\pm$0.003} & 0.928\,{\scriptsize$\pm$0.004} & 0.915\,{\scriptsize$\pm$0.004} & \textbf{0.926}\,{\scriptsize$\pm$0.001} & 0.985\,{\scriptsize$\pm$0.000} \\
  Chen et al.~\cite{chen2019} & 0.909\,{\scriptsize$\pm$0.017} & 0.896\,{\scriptsize$\pm$0.031} & 0.863\,{\scriptsize$\pm$0.005} & 0.871\,{\scriptsize$\pm$0.003} & 0.920\,{\scriptsize$\pm$0.004} & 0.918\,{\scriptsize$\pm$0.002} & 0.921\,{\scriptsize$\pm$0.004} & 0.976\,{\scriptsize$\pm$0.004} \\
  Font et al.~\cite{font2025} & 0.907\,{\scriptsize$\pm$0.019} & 0.886\,{\scriptsize$\pm$0.044} & 0.820\,{\scriptsize$\pm$0.011} & 0.847\,{\scriptsize$\pm$0.003} & 0.886\,{\scriptsize$\pm$0.007} & 0.865\,{\scriptsize$\pm$0.005} & 0.919\,{\scriptsize$\pm$0.003} & 0.976\,{\scriptsize$\pm$0.001} \\
  Gorbett e Blanchard~\cite{gorbett2022} & 0.841\,{\scriptsize$\pm$0.042} & 0.837\,{\scriptsize$\pm$0.078} & 0.880\,{\scriptsize$\pm$0.005} & 0.858\,{\scriptsize$\pm$0.003} & 0.921\,{\scriptsize$\pm$0.006} & 0.922\,{\scriptsize$\pm$0.003} & 0.913\,{\scriptsize$\pm$0.002} & 0.932\,{\scriptsize$\pm$0.031} \\
  ConfidNet~\cite{corbiere2019} & 0.790\,{\scriptsize$\pm$0.032} & 0.757\,{\scriptsize$\pm$0.078} & 0.841\,{\scriptsize$\pm$0.007} & 0.816\,{\scriptsize$\pm$0.003} & 0.901\,{\scriptsize$\pm$0.007} & 0.876\,{\scriptsize$\pm$0.012} & 0.853\,{\scriptsize$\pm$0.006} & 0.792\,{\scriptsize$\pm$0.028} \\
  \midrule
  Fuzzy aditivo, sem vistas (nosso) & 0.914\,{\scriptsize$\pm$0.016} & 0.901\,{\scriptsize$\pm$0.048} & 0.896\,{\scriptsize$\pm$0.004} & 0.875\,{\scriptsize$\pm$0.002} & 0.919\,{\scriptsize$\pm$0.004} & 0.922\,{\scriptsize$\pm$0.004} & 0.925\,{\scriptsize$\pm$0.000} & 0.981\,{\scriptsize$\pm$0.000} \\
  \textbf{Fuzzy aditivo (nosso)} & \textbf{0.916}\,{\scriptsize$\pm$0.012} & \textbf{0.904}\,{\scriptsize$\pm$0.044} & \textbf{0.903}\,{\scriptsize$\pm$0.003} & \textbf{0.886}\,{\scriptsize$\pm$0.002} & \textbf{0.936}\,{\scriptsize$\pm$0.002} & \textbf{0.934}\,{\scriptsize$\pm$0.002} & \textbf{0.926}\,{\scriptsize$\pm$0.002} & 0.982\,{\scriptsize$\pm$0.000} \\
  \bottomrule
\end{tabular}}

\vspace{2pt}
{\footnotesize $^\dagger$ Binário: regressão logística em $P$ e $P^2$; várias classes: idêntico ao softmax.}
\end{table*}

% ==================================================================
% Tabela 3: Erros encontrados
% ==================================================================
\begin{table*}[t]
\centering
\caption{Erros da rede encontrados ao encaminhar ao especialista 10\%, 20\% e 30\% das imagens mais incertas (a cada 1.000 imagens; MNIST a cada 10.000). Entre parênteses, a fração de todos os erros da rede.}
\label{tab:encontrados}
\begin{minipage}[t]{0.495\textwidth}
\centering
\scriptsize\setlength{\tabcolsep}{3pt}
\begin{tabular}[t]{lccc}
  \toprule
  Método & \begin{tabular}[c]{@{}c@{}}Revisando\\10\%\end{tabular} & \begin{tabular}[c]{@{}c@{}}Revisando\\20\%\end{tabular} & \begin{tabular}[c]{@{}c@{}}Revisando\\30\%\end{tabular} \\
  \midrule
  \multicolumn{4}{@{}l}{\textit{HAM10000 (todas)} --- 139 erros} \\
  Ao acaso & 14 (10\%) & 28 (20\%) & 42 (30\%) \\
  Softmax (MSP) & 49 (35\%) & 82 (59\%) & 111 (80\%) \\
  MC Dropout & 48 (34\%) & 82 (59\%) & 111 (80\%) \\
  TTA & 49 (35\%) & 87 (62\%) & 111 (80\%) \\
  Chen et al.~\cite{chen2019} & 61 (44\%) & 104 (75\%) & 122 (88\%) \\
  Gorbett e Blanchard~\cite{gorbett2022} & 54 (39\%) & 88 (64\%) & 110 (79\%) \\
  Só P & \textbf{65 (47\%)} & \textbf{107 (77\%)} & 125 (90\%) \\
  Fuzzy aditivo, sem vistas (nosso) & 63 (45\%) & 106 (76\%) & 125 (90\%) \\
  \textbf{Fuzzy aditivo (nosso)} & 64 (46\%) & 106 (76\%) & \textbf{126 (91\%)} \\
  \midrule
  \multicolumn{4}{@{}l}{\textit{HAM10000 (3000)} --- 169 erros} \\
  Ao acaso & 17 (10\%) & 34 (20\%) & 51 (30\%) \\
  Softmax (MSP) & 34 (20\%) & 82 (48\%) & 126 (74\%) \\
  MC Dropout & 34 (20\%) & 80 (48\%) & 126 (74\%) \\
  TTA & 37 (22\%) & 74 (44\%) & 113 (67\%) \\
  Chen et al.~\cite{chen2019} & 58 (34\%) & 109 (64\%) & 147 (87\%) \\
  Gorbett e Blanchard~\cite{gorbett2022} & 55 (33\%) & 101 (60\%) & 131 (77\%) \\
  Só P & 44 (26\%) & 109 (64\%) & \textbf{148 (87\%)} \\
  Fuzzy aditivo, sem vistas (nosso) & 57 (33\%) & \textbf{119 (70\%)} & \textbf{148 (88\%)} \\
  \textbf{Fuzzy aditivo (nosso)} & \textbf{63 (37\%)} & 118 (70\%) & \textbf{148 (87\%)} \\
  \midrule
  \multicolumn{4}{@{}l}{\textit{CIFAR-10 ruído 30\%} --- 135 erros} \\
  Ao acaso & 13 (10\%) & 27 (20\%) & 40 (30\%) \\
  Softmax (MSP) = Só P & 53 (40\%) & 84 (62\%) & 98 (73\%) \\
  MC Dropout & 49 (36\%) & 78 (58\%) & 94 (70\%) \\
  TTA & 38 (28\%) & 71 (53\%) & 96 (71\%) \\
  Chen et al.~\cite{chen2019} & 53 (39\%) & 88 (66\%) & 108 (80\%) \\
  Gorbett e Blanchard~\cite{gorbett2022} & 56 (41\%) & 93 (69\%) & 113 (84\%) \\
  Fuzzy aditivo, sem vistas (nosso) & \textbf{62 (46\%)} & 96 (72\%) & 116 (86\%) \\
  \textbf{Fuzzy aditivo (nosso)} & \textbf{62 (46\%)} & \textbf{99 (74\%)} & \textbf{119 (88\%)} \\
  \midrule
  \multicolumn{4}{@{}l}{\textit{CIFAR-100} --- 269 erros} \\
  Ao acaso & 27 (10\%) & 54 (20\%) & 81 (30\%) \\
  Softmax (MSP) = Só P & 74 (28\%) & 136 (50\%) & 187 (70\%) \\
  MC Dropout & 77 (29\%) & 138 (51\%) & 186 (69\%) \\
  TTA & 53 (20\%) & 114 (42\%) & 176 (65\%) \\
  Chen et al.~\cite{chen2019} & 76 (28\%) & 137 (51\%) & 187 (69\%) \\
  Gorbett e Blanchard~\cite{gorbett2022} & 74 (28\%) & 136 (50\%) & 184 (68\%) \\
  Fuzzy aditivo, sem vistas (nosso) & 78 (29\%) & 140 (52\%) & 189 (70\%) \\
  \textbf{Fuzzy aditivo (nosso)} & \textbf{81 (30\%)} & \textbf{145 (54\%)} & \textbf{194 (72\%)} \\
  \bottomrule
\end{tabular}
\end{minipage}\hfill
\begin{minipage}[t]{0.495\textwidth}
\centering
\scriptsize\setlength{\tabcolsep}{3pt}
\begin{tabular}[t]{lccc}
  \toprule
  Método & \begin{tabular}[c]{@{}c@{}}Revisando\\10\%\end{tabular} & \begin{tabular}[c]{@{}c@{}}Revisando\\20\%\end{tabular} & \begin{tabular}[c]{@{}c@{}}Revisando\\30\%\end{tabular} \\
  \midrule
  \multicolumn{4}{@{}l}{\textit{CIFAR-10 (ResNet)} --- 71 erros} \\
  Ao acaso & 7 (10\%) & 14 (20\%) & 21 (30\%) \\
  Softmax (MSP) = Só P & 44 (62\%) & 61 (87\%) & 67 (95\%) \\
  MC Dropout & 44 (61\%) & 61 (87\%) & 67 (95\%) \\
  TTA & 45 (63\%) & \textbf{64 (90\%)} & \textbf{68 (96\%)} \\
  Chen et al.~\cite{chen2019} & 42 (60\%) & 60 (84\%) & 67 (94\%) \\
  Gorbett e Blanchard~\cite{gorbett2022} & 43 (61\%) & 60 (85\%) & 66 (93\%) \\
  Fuzzy aditivo, sem vistas (nosso) & 43 (61\%) & 59 (84\%) & 66 (93\%) \\
  \textbf{Fuzzy aditivo (nosso)} & \textbf{48 (68\%)} & \textbf{64 (90\%)} & 67 (95\%) \\
  \midrule
  \multicolumn{4}{@{}l}{\textit{CIFAR-10 (VGG)} --- 63 erros} \\
  Ao acaso & 6 (10\%) & 13 (20\%) & 19 (30\%) \\
  Softmax (MSP) = Só P & 41 (65\%) & 54 (87\%) & 58 (93\%) \\
  MC Dropout & 40 (64\%) & 54 (86\%) & 57 (92\%) \\
  TTA & 43 (69\%) & \textbf{57 (90\%)} & \textbf{60 (96\%)} \\
  Chen et al.~\cite{chen2019} & 37 (59\%) & 54 (86\%) & 59 (94\%) \\
  Gorbett e Blanchard~\cite{gorbett2022} & 39 (62\%) & 54 (86\%) & \textbf{60 (95\%)} \\
  Fuzzy aditivo, sem vistas (nosso) & 40 (65\%) & 54 (86\%) & 59 (94\%) \\
  \textbf{Fuzzy aditivo (nosso)} & \textbf{44 (70\%)} & 56 (90\%) & 59 (95\%) \\
  \midrule
  \multicolumn{4}{@{}l}{\textit{SVHN} --- 46 erros} \\
  Ao acaso & 5 (10\%) & 9 (20\%) & 14 (30\%) \\
  Softmax (MSP) = Só P & \textbf{35 (76\%)} & 41 (90\%) & \textbf{43 (94\%)} \\
  MC Dropout & 32 (69\%) & 41 (89\%) & \textbf{43 (93\%)} \\
  TTA & 30 (66\%) & 41 (89\%) & \textbf{43 (94\%)} \\
  Chen et al.~\cite{chen2019} & 34 (75\%) & 41 (89\%) & \textbf{43 (94\%)} \\
  Gorbett e Blanchard~\cite{gorbett2022} & \textbf{35 (76\%)} & 41 (88\%) & 42 (91\%) \\
  Fuzzy aditivo, sem vistas (nosso) & \textbf{35 (77\%)} & 41 (90\%) & \textbf{43 (94\%)} \\
  \textbf{Fuzzy aditivo (nosso)} & \textbf{35 (76\%)} & \textbf{42 (90\%)} & \textbf{43 (94\%)} \\
  \midrule
  \multicolumn{4}{@{}l}{\textit{MNIST} --- 93 erros (a cada 10.000)} \\
  Ao acaso & 9 (10\%) & 19 (20\%) & 28 (30\%) \\
  Softmax (MSP) = Só P & 92 (99\%) & \textbf{93 (100\%)} & \textbf{93 (100\%)} \\
  MC Dropout & 92 (99\%) & \textbf{93 (100\%)} & \textbf{93 (100\%)} \\
  TTA & \textbf{93 (100\%)} & \textbf{93 (100\%)} & \textbf{93 (100\%)} \\
  Chen et al.~\cite{chen2019} & 91 (97\%) & \textbf{93 (100\%)} & \textbf{93 (100\%)} \\
  Gorbett e Blanchard~\cite{gorbett2022} & 79 (85\%) & 84 (90\%) & 87 (94\%) \\
  Fuzzy aditivo, sem vistas (nosso) & 92 (99\%) & \textbf{93 (100\%)} & \textbf{93 (100\%)} \\
  \textbf{Fuzzy aditivo (nosso)} & \textbf{93 (100\%)} & \textbf{93 (100\%)} & \textbf{93 (100\%)} \\
  \bottomrule
\end{tabular}
\end{minipage}
\end{table*}

% ==================================================================
% Tabela 4: Ganho sobre o softmax
% ==================================================================
\begin{table}[t]
\centering
\caption{Ganho de AUROC do fuzzy aditivo sobre o softmax por cenário.}
\label{tab:ganho}
\resizebox{\columnwidth}{!}{%
\begin{tabular}{lccccc}
  \toprule
  Cenário & Classes & \begin{tabular}[c]{@{}c@{}}Erro da\\rede\end{tabular} & Softmax & \begin{tabular}[c]{@{}c@{}}Fuzzy\\aditivo\end{tabular} & Ganho \\
  \midrule
  MNIST & 10 & 0.9\% & 0.985 & 0.982 & $-$0.003 \\
  SVHN & 10 & 4.6\% & 0.926 & 0.926 & $-$0.001 \\
  CIFAR-10 (ResNet) & 10 & 7.1\% & 0.928 & 0.936 & +0.008 \\
  CIFAR-10 (VGG) & 10 & 6.3\% & 0.915 & 0.934 & +0.019 \\
  CIFAR-100 & 100 & 26.9\% & 0.866 & 0.886 & +0.020 \\
  HAM10000 (todas) & 2 & 13.9\% & 0.865 & 0.916 & +0.052 \\
  HAM10000 (3000) & 2 & 16.9\% & 0.819 & 0.904 & +0.085 \\
  CIFAR-10 ruído 30\% & 10 & 13.5\% & 0.797 & 0.903 & +0.106 \\
  \bottomrule
\end{tabular}}
\end{table}

% ==================================================================
% Tabela 5: Significância
% ==================================================================
\begin{table}[t]
\centering
\caption{Rodadas (sementes) em que o fuzzy aditivo ganha com significância (IC 95\%, bootstrap pareado) de Chen et al. e de Só P.}
\label{tab:signif}
\resizebox{\columnwidth}{!}{%
\begin{tabular}{lcccc}
  \toprule
  Cenário & \begin{tabular}[c]{@{}c@{}}Ganha de\\Chen\end{tabular} & \begin{tabular}[c]{@{}c@{}}Perde para\\Chen\end{tabular} & \begin{tabular}[c]{@{}c@{}}Ganha de Só\\P\end{tabular} & Sementes \\
  \midrule
  HAM10000 (todas) & 1 & 0 & 2 & 5 \\
  HAM10000 (3000) & 1 & 0 & 2 & 3 \\
  CIFAR-10 ruído 30\% & 3 & 0 & 3 & 3 \\
  CIFAR-100 & 3 & 0 & 3 & 3 \\
  CIFAR-10 (ResNet) & 3 & 0 & 3 & 3 \\
  CIFAR-10 (VGG) & 3 & 0 & 3 & 3 \\
  SVHN & 2 & 0 & 0 & 3 \\
  MNIST & 2 & 0 & 0 & 3 \\
  \midrule
  \textbf{Total} & \textbf{18} & \textbf{0} & \textbf{16} & \textbf{26} \\
  \bottomrule
\end{tabular}}
\end{table}


\subsection{Desempenho na detecção}
A Tabela~\ref{tab:auroc} mostra a AUROC. O modelo fuzzy aditivo tem a maior média em seis dos oito cenários e empata com o softmax no
SVHN (0,926); no MNIST, em que a rede erra em menos de 1\% das imagens, fica 0,004 abaixo do melhor detector. Ele também tem o menor E-AURC
nesses mesmos seis cenários, por exemplo 14,8 contra 22,0 de Chen et al.\ e 45,4 do MSP no CIFAR-10 com ruído.
A Tabela~\ref{tab:ganho} mostra que o ganho sobre o softmax cresce quando o softmax é menos confiável: é desprezível no MNIST e no SVHN,
de 0,01 a 0,02 no CIFAR limpo, de 0,05 a 0,09 no HAM10000 e de 0,11 no CIFAR-10 com ruído nos rótulos. Em relação a Chen et al., o concorrente
com treino mais próximo, o modelo fuzzy é significativamente melhor em 18 de 26 rodadas e significativamente pior em nenhuma (Tabela~\ref{tab:signif}).

No HAM10000 o quadro é mais equilibrado. O modelo fuzzy (0,916), Chen et al.\ (0,909) e o modelo só com a saída (0,907) ficam próximos,
e o ganho sobre Chen et al.\ é significativo em apenas uma de cinco sementes. Encaminhando 20\% das lesões, os três métodos encontram
106, 104 e 107 de cada 139 erros por 1.000 lesões (Tabela~\ref{tab:encontrados}), contra 82 do MSP e do MC Dropout. Nesse cenário,
a vantagem do modelo fuzzy é a explicação, não a ordenação. O Mahalanobis sai invertido no HAM10000 (AUROC de 0,18): a maior parte dos erros
são lesões de aparência típica, próximas dos dados de treino, e não lesões incomuns.

A variante de uma passagem, sem vistas, mantém a maior parte do ganho onde ele importa (0,914 no HAM10000 e 0,896 no CIFAR-10 com ruído),
mas no CIFAR-10 limpo as vistas são necessárias para superar o softmax (0,919 contra 0,928).

\subsection{Regras e estabilidade}
As regras são curtas e coerentes com a intuição. No HAM10000 (semente 42), estar no limiar aumenta o log da chance de erro
($p$ \term{médio}: $+1{,}32$), um salto final baixo o diminui ($-1{,}53$) e um desvio do caminho alto entre as vistas o aumenta ($+1{,}16$).
No CIFAR-10 com ruído, as regras dominantes tratam do fim do caminho: concordância final \term{baixa} $+1{,}57$, concordância final
\term{alta} $-2{,}18$. Em dez reamostragens bootstrap, a correlação de Spearman entre os pesos das regras fica entre 0,91 e 0,94 no HAM10000
(todas as lesões) e entre 0,90 e 0,98 nos cenários do CIFAR. A Fig.~\ref{fig:pertinencias} mostra as funções de pertinência e os pesos.

\begin{figure}[t]
\centering
\IfFileExists{figuras/pertinencias.png}{\includegraphics[width=\columnwidth]{figuras/pertinencias.png}}{\fbox{\parbox{0.9\columnwidth}{\centering\vspace{1cm}
coloque figuras/pertinencias.png (gerada por mostrar\_pertinencias.py)\vspace{1cm}}}}
\caption{Funções de pertinência (acima) e pesos das regras (abaixo) do modelo fuzzy aditivo no HAM10000. Um peso positivo aumenta o risco de erro.}
\label{fig:pertinencias}
\end{figure}

\subsection{Verificações com incerteza real}
\textit{Biópsia.} No HAM10000 (todas as lesões, cinco sementes), a rede erra em 29,1\% das lesões confirmadas por histopatologia e em 3,1\%
das confirmadas por acompanhamento, consenso ou microscopia confocal: os casos em que a rede hesita são os mesmos que os médicos também
mandaram para biópsia.
\textit{Outras fotografias da mesma lesão.} O HAM10000 traz fotos extras de muitas lesões do teste, nunca usadas no treino. A rede muda de
resposta em 21\% dessas lesões. Aqui, a distância ao limiar prevê melhor essa mudança (AUROC de 0,72) e o modelo fuzzy é menos eficaz (0,63):
ele foi ajustado para prever erros, não instabilidade entre aquisições.
\textit{Discordância humana.} No CIFAR-10H~\cite{peterson2019}, com cerca de 50 rótulos humanos por imagem de teste, o modelo fuzzy tem a
maior correlação com a discordância humana sob ruído nos rótulos (Spearman de 0,28 contra 0,11 do MSP); nas redes limpas, fica acima do MSP,
mas abaixo de Gorbett e Blanchard (0,33 contra 0,37 com a ResNet-32 e 0,30 contra 0,34 com a VGG-16).

% =====================================================================
\section{Discussão e Limitações}
As características acrescentam mais quando a saída é menos confiável: sob ruído nos rótulos, as camadas finais memorizam rótulos errados e o
softmax fica confiante demais, enquanto as sondas ainda revelam que os blocos anteriores discordam. Em problemas fáceis (MNIST, SVHN), o
softmax já está perto do teto e a leitura interna não acrescenta nada. No HAM10000, a saída binária já é um sinal forte quando se leva o
limiar em conta, então as características internas fornecem sobretudo a explicação.
O detector é treinado com os erros da rede; com poucos erros, como no subconjunto de 3.000 lesões, a variação entre sementes é grande.
As características de vistas multiplicam o custo de inferência por $V$, embora a variante de uma passagem mantenha a maior parte do ganho.
Por fim, todos os detectores foram treinados por validação cruzada no conjunto de teste da mesma rede; se as regras se transferem para outra
arquitetura fica para trabalhos futuros.

% =====================================================================
\section{Conclusão}
Ler uma rede congelada por dentro, pelo caminho da decisão ao longo dos blocos, pela estabilidade entre vistas e pela familiaridade da entrada,
produz uma pontuação de erro tão precisa quanto a dos melhores detectores concorrentes e que, ao contrário deles, pode ser escrita como uma
dúzia de regras fuzzy com pesos. O ganho sobre a confiança da própria rede é maior quando essa confiança é menos confiável, como sob ruído nos
rótulos e na classificação de lesões de pele. Trabalhos futuros vão estudar se as regras se transferem entre arquiteturas, por exemplo com
pertinências fuzzy tipo 2 intervalares que representem a discordância entre redes, e validar a política de encaminhamento com dermatologistas.

% Duplo-cego: sem agradecimentos na submissão.
% \section*{Agradecimentos}

\bibliographystyle{IEEEtran}
\bibliography{refs}

\end{document}
